% !TEX root = ../../ctfp-print.tex

\lettrine[lhang=0.17]{そ}{ろそろ集合について少し話す時が来た。}
数学者は集合論に対して愛憎半ばする関係を持っている。集合論は数学の「アセンブリ言語」だ――少なくとも以前はそうだった。圏論は、ある程度、集合論から距離を置こうとしている。たとえば、すべての集合からなる集合は存在しないことは知られているが、すべての集合からなる圏 $\Set$ は存在する。それは良いことだ。一方で、圏の任意の二対象間の射は集合をなすと仮定している。それをhom集合と呼んでいる。公平を期すために言えば、射が集合をなさない圏論の分野もある。その場合、射は別の圏の対象となる。hom集合ではなくhom対象を用いる圏を\newterm{豊穣圏}（enriched categories）と呼ぶ。しかし以下では、昔ながらのhom集合を持つ圏に留まることにする。

集合は、圏論の対象の外では、最も特徴のない「のっぺりした塊」に近いものだ。集合は元を持つが、それらの元についてあまり多くは言えない。有限集合があれば、元を数えることができる。無限集合の元は基数を使って数えることができる。自然数の集合は実数の集合より小さい――両方とも無限だが。しかし意外なことに、有理数の集合は自然数の集合と同じ大きさだ。

それ以外は、集合に関するすべての情報は、集合の間の関数――特に同型射と呼ばれる可逆な関数――によって符号化できる。あらゆる実用的な目的のために、同型な集合は同一だと見なしてよい。基礎的な数学者の怒りを買う前に、等しさと同型の区別が根本的に重要であることを説明しておこう。実際、これは最新の数学の分野であるホモトピー型理論（Homotopy Type Theory、HoTT）の主要な関心事の一つだ。HoTTに言及するのは、それが計算から着想を得た純粋数学の理論であり、その主要な提唱者の一人であるVladimir Voevodskが、Rocq定理証明器を研究している間に大きな悟りを得たからだ。数学とプログラミングの相互作用は双方向に働く。

集合について重要な教訓は、異なる元からなる集合を比較しても構わないということだ。たとえば、自然変換の集合がある射の集合と同型だと言うことができる。なぜなら集合はただの集合だからだ。この場合の同型とは、一方の集合のすべての自然変換に対して他方の集合のユニークな射が存在し、逆もまた成り立つことを意味する。それらは互いに対応づけられる。異なる圏の対象であればリンゴとオレンジを比較できないが、リンゴの集合とオレンジの集合は比較できる。圏論的問題を集合論的問題に変換することで、必要な洞察が得られたり、価値ある定理を証明できることも多い。

\section{Hom関手}

すべての圏には、$\Set$ への写像の標準的な族が備わっている。これらの写像は実は関手であり、圏の構造を保存する。そのような写像を一つ構成しよう。

$\cat{C}$ の中の対象 $a$ を一つ固定し、同じく $\cat{C}$ の中の別の対象 $x$ を選ぶ。hom集合 $\cat{C}(a, x)$ は集合であり、$\Set$ の対象だ。$a$ を固定したまま $x$ を変えると、$\cat{C}(a, x)$ も $\Set$ の中で変化する。したがって、$x$ から $\Set$ への写像が得られる。

\begin{figure}[H]
  \centering
  \includegraphics[width=0.45\textwidth]{images/hom-set.jpg}
\end{figure}

\noindent
hom集合を第二引数の写像として考えていることを強調したい場合、引数の場所として「$-$」を使い、$\cat{C}(a, -)$ という記法を用いる。

対象のこの写像は、射の写像へと容易に拡張される。$\cat{C}$ の任意の二対象 $x$ と $y$ の間の射 $f$ を取ろう。先ほど定義した写像の下で、対象 $x$ は集合 $\cat{C}(a, x)$ に写され、対象 $y$ は $\cat{C}(a, y)$ に写される。この写像が関手であるためには、$f$ は二つの集合の間の関数に写されなければならない：$\cat{C}(a, x) \to \cat{C}(a, y)$

この関数を点ごとに定義しよう。すなわち、各引数について別々に定義する。引数として、$\cat{C}(a, x)$ の任意の元を選ぶ――それを $h$ と呼ぼう。射は端が合えば合成できる。$h$ の終域が $f$ の始域と一致するので、それらの合成：
\[f \circ h \Colon a \to y\]
は $a$ から $y$ への射だ。したがってそれは $\cat{C}(a, y)$ の元だ。

\begin{figure}[H]
  \centering
  \includegraphics[width=0.45\textwidth]{images/hom-functor.jpg}
\end{figure}

\noindent
こうして $\cat{C}(a, x)$ から $\cat{C}(a, y)$ への関数が見つかった。これが $f$ の像として機能する。混乱の恐れがなければ、この持ち上げられた関数を $\cat{C}(a, f)$ と書き、射 $h$ への作用を：
\[\cat{C}(a, f) h = f \circ h\]
と表す。この構成はあらゆる圏で機能するので、Haskellの型の圏でも機能するはずだ。Haskellでは、hom関手は\code{Reader}関手としてよく知られている：

\src{snippet01}

\src{snippet02}
では、hom集合の始域を固定する代わりに、終域を固定したらどうなるかを考えよう。言い換えれば、写像 $\cat{C}(-, a)$ も関手かどうかという問いだ。そうなのだが、共変ではなく反変だ。なぜなら、同じように射の端を合わせると、$\cat{C}(a, -)$ の場合の前合成ではなく、$f$ による後合成になるからだ。

この反変関手はHaskellですでに見た。\code{Op} と呼んだものだ：

\src{snippet03}

\src{snippet04}
最後に、両方の対象を変化させると、プロ関手 $\cat{C}(-, =)$ が得られる。これは第一引数で反変、第二引数で共変だ（二つの引数が独立して変化しうることを強調するために、第二のプレースホルダーとして二重のダッシュを使う）。このプロ関手は関手性について話したときにすでに見た：

\src{snippet05}
重要な教訓は、この観察はどんな圏でも成り立つということだ。対象からhom集合への写像は関手的だ。反変性が反対圏からの写像と等価なので、この事実を簡潔に次のように述べられる：
\[C(-, =) \Colon \cat{C}^\mathit{op} \times \cat{C} \to \Set\]

\section{表現可能関手}

$\cat{C}$ の対象 $a$ を選ぶごとに、$\cat{C}$ から $\Set$ への関手が得られることを見た。このような $\Set$ への構造保存写像はしばしば\newterm{表現}（representation）と呼ばれる。$\cat{C}$ の対象と射を、$\Set$ の集合と関数として表現しているわけだ。

関手 $\cat{C}(a, -)$ 自体を表現可能（representable）と呼ぶことがある。より一般的には、ある $a$ の選択に対してhom関手と自然同型な任意の関手 $F$ を\newterm{表現可能}と呼ぶ。そのような関手は必然的に $\Set$ 値でなければならない。$\cat{C}(a, -)$ がそうであるからだ。

以前、同型な集合を同一だと見なすことが多いと述べた。より一般的には、圏の中の同型な\emph{対象}を同一と見なす。なぜなら、対象には射を通じた他の対象（および自身）との関係以外の構造がないからだ。

たとえば、最初に集合でモデル化したモノイドの圏 $\cat{Mon}$ について話した。しかし、射としてはそれらの集合のモノイド構造を保存する関数だけを選ぶよう注意した。だから $\cat{Mon}$ の二つの対象が同型、すなわち可逆な射が存在するならば、それらはまったく同じ構造を持つ。それらが基づく集合と関数を調べれば、一方のモノイドの単位元が他方の単位元に写され、二つの元の積がそれらの像の積に写されることがわかるだろう。

同じ論理は関手にも適用できる。二つの圏の間の関手は、自然変換が射の役割を担う圏をなす。したがって、二つの関手の間に可逆な自然変換が存在するならば、それらは同型であり、同一と見なすことができる。

この観点から表現可能関手の定義を分析してみよう。$F$ が表現可能であるためには：$\cat{C}$ の中の対象 $a$ が存在すること；$\cat{C}(a, -)$ から $F$ への自然変換 $\alpha$ が一つ存在すること；逆方向の自然変換 $\beta$ が一つ存在すること；そしてそれらの合成が恒等自然変換であることが必要だ。

ある対象 $x$ における $\alpha$ の成分を見てみよう。それは $\Set$ の関数だ：
\[\alpha_x \Colon \cat{C}(a, x) \to F x\]
この変換の自然条件は、$x$ から $y$ への任意の射 $f$ に対して、次の図式が可換であることだ：
\[F f \circ \alpha_x = \alpha_y \circ \cat{C}(a, f)\]
Haskellでは、自然変換を多相関数で置き換える：

\src{snippet06}
オプションの \code{forall} 量化子付きで。自然条件

\src{snippet07}
はパラメトリシティによって自動的に満たされる（これは以前述べた「タダで得られる定理」の一つだ）。左辺の \code{fmap} は関手 $F$ によって定義され、右辺の \code{fmap} はreader関手によって定義されることを理解したうえで。readerの \code{fmap} は単なる関数の前合成なので、より明示的にできる。$\cat{C}(a, x)$ の元である $h$ に作用させると、自然条件は次のように簡略化される：

\src{snippet08}
もう一方の変換 \code{beta} は逆方向に進む：

\src{snippet09}
これは自然条件を満たさなければならず、\code{alpha} の逆でなければならない：

\begin{snip}{text}
alpha . beta = id = beta . alpha
\end{snip}
後で、$\cat{C}(a, -)$ から任意の $\Set$ 値関手への自然変換は、$F a$ が空でない限り常に存在することがわかる（米田の補題）が、それが必ずしも可逆とは限らない。

Haskellでリスト関手と \code{a} として \code{Int} を使った例を示そう。次の自然変換がその役割を果たす：

\src{snippet10}
任意に数 12 を選び、それで単要素リストを作った。そのリストに関数 \code{h} を \code{fmap} することで、\code{h} が返す型のリストが得られる。（実際には整数のリストと同数のそのような変換が存在する。）

自然条件は \code{map}（\code{fmap} のリスト版）の合成可能性と等価だ：

\src{snippet11}
しかし逆変換を見つけようとすれば、任意の型 \code{x} のリストから \code{x} を返す関数へ進む必要がある：

\src{snippet12}
たとえば \code{head} を使ってリストから \code{x} を取り出すことを考えるかもしれないが、空リストではうまくいかない。ここでうまく機能する型 \code{a}（\code{Int} の代わりに）はないことに注意されたい。したがって、リスト関手は表現可能ではない。

Haskellの（内）関手がコンテナに少し似ていると話したことを覚えているだろうか？同様に、表現可能関手を関数呼び出しのメモ化された結果を格納するコンテナとして考えることができる（Haskellのhom集合の元は単なる関数だ）。表現対象、すなわち $\cat{C}(a, -)$ の型 $a$ は、表形式の関数値にアクセスするためのキーの型と考えられる。\code{alpha} と呼んでいた変換は \code{tabulate} と呼ばれ、その逆 \code{beta} は \code{index} と呼ばれる。以下が（少し簡略化した）\code{Representable} クラスの定義だ：

\src{snippet13}
表現する型（ここでは \code{Rep f} と呼ばれる私たちの $a$）が \code{Representable} の定義の一部であることに注意されたい。アスタリスクは \code{Rep f} が型（型コンストラクタやその他のより奇妙な種ではなく）であることを意味する。

無限リスト、つまりストリームは空になり得ないため、表現可能だ。

\src{snippet14}
これらを \code{Integer} を引数として取る関数のメモ化された値として考えることができる。（厳密には、非負の自然数を使うべきだが、コードを複雑にしたくなかった。）

そのような関数を \code{tabulate} するには、値の無限ストリームを作る。もちろん、これはHaskellが遅延評価であるためにのみ可能だ。値は必要に応じて評価される。メモ化された値へのアクセスには \code{index} を使う：

\src{snippet15}
任意の戻り値型を持つ関数の族全体をカバーする単一のメモ化スキームを実装できるのは興味深い。

反変関手の表現可能性も同様に定義されるが、$\cat{C}(-, a)$ の第二引数を固定する。あるいは同等に、$\cat{C}^\mathit{op}(a, -)$ は $\cat{C}(-, a)$ と同じなので、$\cat{C}^\mathit{op}$ から $\Set$ への関手を考えてもよい。

表現可能性には興味深いひねりがある。デカルト閉圏では、hom集合は内部的に指数対象として扱えることを思い出そう。hom集合 $\cat{C}(a, x)$ は $x^a$ と等価であり、表現可能関手 $F$ に対して $-^a = F$ と書ける。

面白半分に両辺の対数を取ってみると：

\[a = \log F\]

もちろんこれは純粋に形式的な変換だが、対数の性質をいくつか知っていれば、非常に役に立つ。特に、積型に基づく関手は和型で表現できること、そして和型の関手は一般に表現可能ではない（例：リスト関手）ことがわかる。

最後に、表現可能関手は同じものの二つの異なる実装を与えてくれることに注意しよう――一方は関数、もう一方はデータ構造だ。それらは完全に同じ内容を持っている――同じキーを使って同じ値が取り出される。それが私が言っていた「同じさ」の意味だ。自然同型な二つの関手は、その内容に関する限り同一だ。一方で、二つの表現は異なる方法で実装されることが多く、異なるパフォーマンス特性を持つことがある。メモ化はパフォーマンス向上のために使われ、実行時間を大幅に短縮できる。同じ基礎的計算の異なる表現を生成できることは、実践的に非常に価値がある。だから意外にも、パフォーマンスをまったく気にしないにもかかわらず、圏論は実用的価値を持つ代替実装を探索する豊富な機会を提供してくれる。

\section{練習問題}

\begin{enumerate}
  \tightlist
  \item
        Hom関手が $\cat{C}$ の恒等射を $\Set$ の対応する恒等関数に写すことを示せ。
  \item
        \code{Maybe} が表現可能でないことを示せ。
  \item
        \code{Reader} 関手は表現可能か？
  \item
        \code{Stream} 表現を使って、引数を二乗する関数をメモ化せよ。
  \item
        \code{Stream} に対する \code{tabulate} と \code{index} が互いの逆であることを示せ。（ヒント：帰納法を使え。）
  \item
        関手：

        \begin{snip}{haskell}
Pair a = Pair a a
\end{snip}
        は表現可能だ。それを表現する型を推測できるか？\code{tabulate} と \code{index} を実装せよ。
\end{enumerate}

\section{参考文献}

\begin{enumerate}
  \tightlist
  \item
        \urlref{https://www.youtube.com/watch?v=4QgjKUzyrhM}{表現可能関手}に関するCatstersのビデオ。
\end{enumerate}
